\begin{abstract}
Spectral adaptation---tuning the frequency parameters of a continuous
representation, allocating capacity per band, or injecting a spectral
prior---gives large gains on some reconstruction tasks and almost none on
others. Arbitrary-scale super-resolution (ASISR) is the clearest case of the
latter: three years of decoder design on top of LIIF have moved the standard
benchmark by roughly a third of a decibel, and a recent controlled study
reports architectural saturation. We ask whether one can tell \emph{in
advance}, without training and without ground truth, whether spectral
adaptation will help on a given task and frequency band, how much it can
help, and what it costs when the spectral prior is wrong.
We (i) show that the common forms of spectral adaptation amount to choosing a
prior spectrum (Bochner); (ii) compute a per-band linear recoverability
$\Rlin$ from the observation operator alone, which upper-bounds every
stationary-kernel linear estimator; and (iii) derive a closed-form mismatch
cost $\Delta(R,\kappa)$ whose asymmetry splits bands into three regimes:
wide/low-bias representations are safe where $R$ is high, the spectrum must
be matched in between, and only conservative priors or learned (non-linear)
priors are safe where $R$ is low.
Pre-registered experiments on single-image ASISR at $\times2$--$\times8$,
sparse-pixel reconstruction, multi-frame super-resolution and 1D EEG confirm
the predicted mismatch-cost shape and asymmetry, locate the decoder
bottleneck of ASISR in the low-$R$ regime where four decoders (LIIF, LTE,
SC-INR and a variant) differ by at most 0.03 in band recovery, and locate the
benefit of the cell/footprint readout (LTE and SC-INR, seed 1) in the bands that the trained network, not the linear
budget, partially recovers. We report the pre-registered
criteria that failed alongside those that held.
\end{abstract}
